\chapter{Адаптивно обучение: теоретичен апарат}

Настоящата глава представя основния теоретичен апарат на адаптивното обучение в духа на подхода, развит от \textcite{GeorgeEvans_2001}. Целта е да се формализира процесът, чрез който икономическите агенти формират и актуализират очакванията си в условия на непълна информация относно структурната динамика на икономиката.


\section{Рационални очаквания}

При парадигмата на рационалните очаквания се приема, че икономическите агенти познават истинската структура на икономиката и използват цялата налична информация ефективно при формирането на своите прогнози.

Въпреки че хипотезата за рационални очаквания предлага елегантна теоретична рамка, тя предполага значителна когнитивна способност от страна на икономическите агенти. Адаптивното обучение предоставя алтернативен подход, при който агентите постепенно актуализират своите вярвания въз основа на наблюдаваните данни, като не се изисква предварително познаване на структурата на икономиката.

\section{Възприета динамика (PLM)}

В теорията на адаптивното обучение се приема, че икономическите агенти не разполагат с пълна информация относно структурната динамика на икономиката. Вместо това те формират прогнозите си въз основа на наблюдаваните данни, използвайки иконометричен регресионен модел, който
свързва ендогенните променливи с наблюдавани регресори. Този модел се означава като \emph{възприета динамика} или \emph{Perceived Law of Motion} (PLM).

За да илюстрираме механизма на формиране на прогнозите, в по-нататъшното изложение разглеждаме опростен линеен пример на PLM с една ендогенна променлива.

\begin{equation}
y_{t+1} = a + B^\top z_{t} + \varepsilon_{t+1},
\label{eq:plm}
\end{equation}
където $a \in \mathbb{R}$ е скаларен параметър,
$B \in \mathbb{R}^{m}$ е вектор от коефициенти
$z_{t} \in \mathbb{R}^m$ е вектор от наблюдавани екзогенни или лагнати
променливи, налични към момента $t$, а $\varepsilon_{t+1}$ е стохастичен шум, отразяващ прогнозните грешки в рамките на възприетия от агентите модел.

На базата на възприетата динамика (PLM) агентите формират прогноза за бъдещата стойност на ендогенната променлива, базирана на информацията, налична към момента $t$.
\begin{equation}
y_{t+1}^e = a_{t} + B_{t}^\top z_{t},
\label{eq:plm_expectation}
\end{equation}
където $(a_{t}, B_{t})$ са текущите оценки на параметрите на PLM.

\section{Рекурсивен метод на най-малките квадрати  (RLS) като алгоритъм за обучение}

\begin{equation}
x_{t} \equiv
\begin{pmatrix}
1 \\
z_{t}
\end{pmatrix},
\qquad
\phi_t \equiv
\begin{pmatrix}
a_t \\
B_t
\end{pmatrix}.
\end{equation}
Нека $k \equiv 1+m$. Тогава $x_{t} \in \mathbb{R}^k$ и $\phi_t \in \mathbb{R}^k$, а уравнението на PLM може да бъде записано като
\begin{equation}
y_{t+1} = \phi_t^\top x_{t} + \varepsilon_{t+1}.
\end{equation}

При постъпване на ново наблюдение $(x_{t}, y_{t+1})$ агентите трябва да актуализират своите убеждения относно вектора от параметри $\phi$. Вместо да преизчисляват оценката от нулата на база цялата налична история от данни, те използват рекурсивен алгоритъм, който обновява предишната оценка ${\phi}_t$ с новата информация.

Оценката на параметрите на възприетата динамика (PLM) се осъществява чрез алгоритъма \emph{Recursive Least Squares} (RLS). RLS представлява рекурсивна формулировка на стандартния метод на най-малките квадрати (МНМК) и позволява последователно обновяване на оценките на параметрите при постъпване на нови наблюдения. RLS e стандартен алгоритърм в литературата по адаптивно обучение. Подробното му извеждане е представено в \textcite[с. 17-20]{ljung1983theory}.

Стандартната оценка по МНМК в момент $t$ е:
\begin{equation}
{\phi}_t = \left(\sum_{s=1}^{t} x_{s} x_{s}^\top\right)^{-1} \sum_{s=1}^{t} x_{s} y_{s+1}.
\end{equation}
Дефинираме нормализираните моменти:
\begin{equation}
R_t = \frac{1}{t}\sum_{s=1}^{t} x_{s} x_{s}^\top, \qquad
p_t = \frac{1}{t}\sum_{s=1}^{t} x_{s} y_{s+1},
\end{equation}
така че ${\phi}_t = R_t^{-1} p_t$.

За нормализирания момент от втори ред имаме:
\begin{align}
R_t &= \frac{1}{t}\sum_{s=1}^{t} x_{s} x_{s}^\top = \frac{1}{t}\left(\sum_{s=1}^{t-1} x_{s} x_{s}^\top + x_{t} x_{t}^\top\right) \notag \\
&= \frac{1}{t}\left((t-1)R_{t-1} + x_{t} x_{t}^\top\right) = \frac{t-1}{t}R_{t-1} + \frac{1}{t}x_{t} x_{t}^\top \notag \\
&= R_{t-1} + \frac{1}{t}(x_{t} x_{t}^\top - R_{t-1}).
\end{align}

Аналогично, за нормализирания момент в числителя:
\begin{align}
p_t &= \frac{1}{t}\sum_{s=1}^{t} x_{s} y_{s+1} = \frac{1}{t}\left(\sum_{s=1}^{t-1} x_{s} y_{s+1} + x_{t} y_{t+1}\right) \notag \\
&= \frac{1}{t}\left((t-1)p_{t-1} + x_{t} y_{t+1}\right) = \frac{t-1}{t}p_{t-1} + \frac{1}{t}x_{t} y_{t+1} \notag \\
&= p_{t-1} + \frac{1}{t}(x_{t} y_{t+1} - p_{t-1}),
\end{align}
където използваме факта, че $p_{t-1} = R_{t-1} {\phi}_{t-1}$ по дефиниция на ${\phi}_{t-1}$.

За да получим рекурсивна форма, изчисляваме $p_t - R_t \phi_{t-1}$:
\begin{align}
p_t - R_t\phi_{t-1} &= \frac{t-1}{t}R_{t-1}\phi_{t-1} + \frac{1}{t}x_t y_{t+1} - \left(\frac{t-1}{t}R_{t-1} + \frac{1}{t}x_t x_t^\top\right)\phi_{t-1} \notag \\
&= \frac{t-1}{t}R_{t-1}\phi_{t-1} + \frac{1}{t}x_t y_{t+1} - \frac{t-1}{t}R_{t-1}\phi_{t-1} - \frac{1}{t}x_t x_t^\top\phi_{t-1} \notag \\
&= \frac{1}{t}x_t(y_{t+1} - x_t^\top\phi_{t-1}).
\end{align}

Тъй като $\phi_t = R_t^{-1} p_t$, получаваме:
\begin{align}
\phi_t - \phi_{t-1} &= R_t^{-1}p_t - \phi_{t-1} = R_t^{-1}(p_t - R_t\phi_{t-1}) \notag \\
&= \frac{1}{t}R_t^{-1} x_t(y_{t+1} - x_t^\top\phi_{t-1}).
\end{align}

Следователно рекурсивната формула е:
\begin{equation}
{\phi}_t = {\phi}_{t-1} + \frac{1}{t}R_t^{-1} x_{t} (y_{t+1} - x_{t}^\top {\phi}_{t-1}),
\label{eq:rls_basic}
\end{equation}
където $(y_{t+1} - x_{t}^\top {\phi}_{t-1})$ е прогнозната грешка.

Формулата \eqref{eq:rls_basic} показва, че при класическия RLS коефициентът на обучение е естествено $\gamma_t = 1/t$. Това съответства на еднакво тегло на всички наблюдения в извадката.

Адаптивното обучение обобщава този подход, като замества специалния коефициент $1/t$ с произволна последователност $\gamma_t$. Рекурсивните правила стават:
\begin{align}
R_t &= R_{t-1} + \gamma_t(x_t x_t^\top - R_{t-1}), \label{eq:R_general}\\
p_t &= p_{t-1} + \gamma_t(x_t y_{t+1} - p_{t-1}). \label{eq:p_general}
\end{align}

Запазваме релацията $\phi_t = R_t^{-1} p_t$. Използвайки $p_{t-1} = R_{t-1}\phi_{t-1}$, може да се покаже (аналогично на извеждането за $\gamma_t = 1/t$), че:
\begin{equation}
p_t - R_t\phi_{t-1} = \gamma_t x_t(y_{t+1} - x_t^\top\phi_{t-1}).
\end{equation}

Следователно рекурсивните формули за адаптивно обучение са:
\begin{align}
\phi_t &= \phi_{t-1} + \gamma_t R_t^{-1} x_t (y_{t+1} - x_t^\top \phi_{t-1}),
\label{eq:RLS_phi} \\
R_t &= R_{t-1} + \gamma_t (x_t x_t^\top - R_{t-1}).
\label{eq:RLS_R}
\end{align}

Това е формулировката, използвана от \textcite[стр. 33-34]{GeorgeEvans_2001}.

В литературата по адаптивно обучение се използват два основни типа коефициенти на обучение \parencite[раздел 15.2]{GeorgeEvans_2001}.

\begin{itemize}
\item \emph{Намаляващ коефициент} (decreasing gain) от вида $\gamma_t = t^{-\beta}$ с $0 < \beta \leq 1$, който осигурява сходимост към равновесие с рационални очаквания при правилно специфициран модел на възприеманата динамика.

\item \emph{Константен коефициент} (constant gain) $\gamma_t = \bar{\gamma}$ с $0 < \bar{\gamma} \leq 1$, който запазва постоянна адаптивност и позволява ефективно следене на структурни промени в икономиката, но не води до сходимост към фиксирана точка.
\end{itemize}

За да се гарантира сходимост при намаляващ коефициент, той трябва да удовлетворява условията $\sum_{t=1}^{\infty} \gamma_t = \infty$ и $\sum_{t=1}^{\infty} \gamma_t^2 < \infty$. 

Първото условие $\sum_{t=1}^{\infty} \gamma_t = \infty$ изисква сумата на коефициентите на обучение да бъде неограничена. Ако $\sum_{t=1}^{\infty} \gamma_t < \infty$, то след достатъчно време сумата на всички бъдещи актуализации ще бъде ограничена, което означава, че ако алгоритъмът е конвергирал към грешна стойност, той няма да разполага с достатъчна обучаваща сила за корекция към правилното равновесие. Второто условие $\sum_{t=1}^{\infty} \gamma_t^2 < \infty$ контролира дисперсията на оценките в стохастична среда. Вариацията на $\phi_t$ зависи от $\sum \gamma_t^2$: ако тази сума е неограничена, стохастичните шокове продължат да влияят на оценката и няма да се получи сходимост към константа. Условието осигурява, че влиянието на шума затихва асимптотично \parencite[с. 124]{GeorgeEvans_2001}.

За коефициенти от вида $\gamma_t = t^{-\beta}$ тези условия са изпълнени при $\beta \in (1/2, 1]$. Стойността $\beta = 1$ съответства на класическия рекурсивен метод на най-малките квадрати, при който всички наблюдения имат еднакво тегло. При $\beta \in (1/2, 1)$ алгоритъмът запазва по-висока адаптивност, тъй като коефициентът на обучение намалява по-бавно и ефективната памет на алгоритъма расте със скорост $t^\beta$ вместо $t$.

В настоящото изложение използваме класическия избор $\gamma_t = t^{-1}$.

\section{Действителна Динамика (ALM)}

Действителната динамика на икономическите променливи се определя от структурния модел на икономиката, в който очакванията на агентите играят активна роля. Като отправна точка разглеждаме \emph{действителната} структурна динамика на икономиката (Actual Law of Motion, ALM), която може да бъде представена в общ вид като:

\begin{equation}
\mathbf{y}_t = F\bigl( \mathbf{y}_{t+1}^e,\; \mathbf{z}_{t},\; \boldsymbol{\eta}_t \bigr),
\label{eq:structural_model}
\end{equation}
където $\mathbf{y}_t \in \mathbb{R}^n$ е вектор от ендогенни променливи, $\mathbf{z}_{t} \in \mathbb{R}^m$ е вектор от наблюдавани екзогенни или лагнати променливи, а $\boldsymbol{\eta}_t \in \mathbb{R}^n$ е вектор от структурни шокове. С $\mathbf{y}_{t+1}^e$ обозначаваме вектора с прогнозите на агентите за бъдещите стойности на ендогенните променливи $\mathbf{y}_{t+1}$, формиран на базата на информацията, налична към момента $t$.

Формулировката \eqref{eq:structural_model} важи и за произволни очаквания на агентите - не само за рационални. Очакванията $\mathbf{y}_{t+1}^e$ могат да бъдат формирани по различни начини: чрез адаптивно учене, чрез прости адаптивни правила, или при хипотезата за рационални очаквания.

За да се илюстрира механизма, чрез който възприетата динамика поражда
действителна динамика, в по-нататъшното изложение се разглежда опростен
скаларен пример, в който както възприетата, така и действителната
динамика са линейни.

В този линеен пример действителната динамика има вида
\begin{equation}
y_t
=
\alpha\, y_{t+1}^e
+
\beta^\top z_{t}
+
\eta_t,
\label{eq:ALM_linear}
\end{equation}
където $\alpha \in \mathbb{R}$ и $\beta \in \mathbb{R}^{m}$ са
структурни параметри, а $\eta_t$ е стохастичен структурен шок.

Замествайки прогнозата $y_{t+1}^e$, генерирана от възприетата динамика (PLM) чрез уравнение \eqref{eq:plm_expectation},
в уравнение \eqref{eq:ALM_linear}, получаваме
\[
y_t
=
\alpha\bigl(a_{t} + B_{t}^\top z_{t}\bigr)
+
\beta^\top z_{t}
+
\eta_t,
\]
или еквивалентно
\begin{equation}
y_t
=
\alpha a_t
+
(\alpha B_t + \beta)^\top z_t
+
\eta_t.
\label{eq:ALM_final}
\end{equation}

Уравнение \eqref{eq:ALM_final} описва реализираната действителна динамика,
която възниква в резултат на използването на прогнозите, формирани чрез
възприетата динамика, в структурния модел на икономиката.


\section{T-mapping}

Със заместването на очакването за бъдещите стойности на ендогенните променливи с тези породени от възприятията на агентите се демонстрира как параметрите на PLM еднозначно определят реализираната форма на действителната динамика.

В рамките на тази формулировка възниква следният централен въпрос на адаптивното обучение: ако агентите продължават да актуализират вектора на текущите оценки на коефициентите от PLM, 
\[
\phi_t \equiv 
\begin{pmatrix} a_t \\ B_t 
\end{pmatrix}. 
\] 
ще се получи ли сходимост на тези вярвания към такава стойност, при която възприеманата и действителната динамика да съвпаднат?

В този контекст се дефинира операторът $T$, който за всяка стойност от вектора $\phi_t$ съпоставя параметрите на породената от него действителна динамика. За разглеждания линеен пример операторът $T$ има вида

\begin{equation} T:
\begin{pmatrix} a\\ B \end{pmatrix} \mapsto
\begin{pmatrix} \alpha a\\ \alpha B + \beta \end{pmatrix}. 
\label{eq:T_mapping} 
\end{equation} 

\section{Равновесие с рационални очаквания}

Неподвижна точка на оператора $T$, обозначена с $\phi^{\ast}$, се нарича \emph{равновесие с рационални очаквания} (Rational Expectations Equilibrium, REE). При това равновесие възприеманата от агентите динамика (PLM) съвпада с действителната динамика (ALM), породена от структурния модел. В икономически смисъл това означава, че агентите са научили истинския модел на развитие на икономиката -- техните вярвания за параметрите на икономическата динамика съответстват на действителните параметри.

Това условие се изразява чрез:

\begin{equation}
\phi^{\ast} = T(\phi^{\ast}),
\label{eq:T_fixed_point}
\end{equation}

\section{E-устойчивост}

Казваме, че дадено равновесие с рационални очаквания $\phi^*$ е 
\emph{E-устойчиво}, ако то е локално асимптотично устойчиво при 
диференциалното уравнение
\begin{equation}
\frac{d\phi}{d\tau} = T(\phi) - \phi,
\label{eq:E_dynamic}
\end{equation}
където $\tau$ е параметър на „времето на обучението“ \parencite[с. 31]{GeorgeEvans_2001}. Дясната страна $T(\phi) - \phi$ изразява разликата между параметрите $T(\phi)$, които биха се породили от действителната динамика при възприемани параметри $\phi$, и самите възприемани параметри. Уравнението описва хипотетичен процес, при който възприеманите параметри се движат в посока на породените от тях, и $\phi^*$ е E-устойчиво, ако малки отклонения от него затихват при този процес. Така E-устойчивостта осигурява критерий за сходимост на адаптивното обучение към равновесието с рационални очаквания.

Връзката между E-устойчивостта и сходимостта на действителния процес на обучение се установява чрез теорията на стохастичната апроксимация. Рекурсивният алгоритъм \eqref{eq:RLS_phi}--\eqref{eq:RLS_R} е дискретен и стохастичен, но тъй като $\gamma_t = t^{-1}$ клони към нула, последователните корекции стават все по-малки, а стохастичните им компоненти се усредняват; в резултат асимптотичното поведение на $\phi_t$ се описва от траекторията на \eqref{eq:E_dynamic} \parencite[с. 127--128]{GeorgeEvans_2001}. Основният резултат на този подход е, че $\phi_t$ може да се сходи единствено към локално асимптотично устойчивите равновесни точки на \eqref{eq:E_dynamic}; сходимост към неустойчиви точки е изключена \parencite[с. 35]{GeorgeEvans_2001}.

\subsection*{Локален анализ}

Нека $\phi^\ast$ е неподвижна точка на $T$. За изследване на локалното поведение се разглежда отклонението
\[
\tilde{\phi}=\phi-\phi^\ast.
\]

За да получим линеаризираната форма на динамиката \eqref{eq:E_dynamic}, изразяваме $\phi = \phi^* + \tilde{\phi}$ и заместваме в \eqref{eq:E_dynamic}:
\begin{equation}
\frac{d(\phi^* + \tilde{\phi})}{d\tau} = T(\phi^* + \tilde{\phi}) - (\phi^* + \tilde{\phi}).
\label{eq:substitution}
\end{equation}

Тъй като $\phi^*$ е константа, $\frac{d\phi^*}{d\tau} = 0$, и получаваме:
\begin{equation}
\frac{d\tilde{\phi}}{d\tau} = T(\phi^* + \tilde{\phi}) - \phi^* - \tilde{\phi}.
\label{eq:deviation_ode}
\end{equation}

Разлагаме $T(\phi^* + \tilde{\phi})$ в ред на Тейлър от първи ред около точката $\phi^*$:
\begin{equation}
T(\phi^* + \tilde{\phi}) = T(\phi^*) + DT(\phi^*)\tilde{\phi} + O(|\tilde{\phi}|^2),
\label{eq:taylor_T}
\end{equation}
където $DT(\phi^*)$ е матрицата на първите частни производни (Якобиан) на $T$, оценена във $\phi^*$, а $O(|\tilde{\phi}|^2)$ обозначава членове от втори и по-висок ред.

Заместваме \eqref{eq:taylor_T} в \eqref{eq:deviation_ode}:
\begin{equation}
\frac{d\tilde{\phi}}{d\tau} = T(\phi^*) + DT(\phi^*)\tilde{\phi} + O(|\tilde{\phi}|^2) - \phi^* - \tilde{\phi}.
\label{eq:substituted}
\end{equation}

След опростяване:
\begin{equation}
\frac{d\tilde{\phi}}{d\tau} = DT(\phi^*)\tilde{\phi} - \tilde{\phi} + O(|\tilde{\phi}|^2) = [DT(\phi^*) - I]\tilde{\phi} + O(|\tilde{\phi}|^2).
\label{eq:almost_linear}
\end{equation}

За достатъчно малки стойности на $|\tilde{\phi}|$, членовете от по-висок ред стават пренебрежимо малки и могат да бъдат изключени. Това дава линеаризираната система:
\begin{equation}
\dot{\tilde{\phi}} = [DT(\phi^\ast)-I]\tilde{\phi},
\label{eq:E_linearized_final}
\end{equation}

Преминаването от анализа на нелинейната динамика \eqref{eq:E_dynamic} 
към линеаризираната система \eqref{eq:E_linearized_final} се обосновава 
от резултат на теорията на обикновените диференциални уравнения за 
локалната устойчивост на нелинейни системи в околност на неподвижна точка, 
чиито Якобиан няма собствени стойности с нулева реална част (т.е. точката 
е хиперболична): в такъв случай локалната устойчивост на нелинейната 
система съвпада с тази на нейната линеаризация \parencite[Proposition 5.6, 
с. 96]{GeorgeEvans_2001}.

Следователно локалната асимптотична устойчивост на $\phi^\ast$ при 
\eqref{eq:E_dynamic} е еквивалентна на тази на $\tilde{\phi}=0$ при 
\eqref{eq:E_linearized_final}, която се определя от собствените стойности 
на $DT(\phi^\ast)-I$: равновесието е асимптотично устойчиво, ако всички 
те имат отрицателни реални части 
\parencite[Proposition 5.6, с. 96]{GeorgeEvans_2001}. Това условие 
гарантира, че малките отклонения на вярванията от $\phi^\ast$ клонят 
към нула при $\tau \to \infty$.

\subsection*{Приложение към линейния пример}

В разглеждания линеен модел операторът $T$ има вида

\[
T(a,B)=
\begin{pmatrix}
\alpha a \\
\alpha B + \beta
\end{pmatrix}.
\]

Якобианът на $T$ е

\[
DT(\phi)=
\begin{pmatrix}
\alpha & 0 \\
0 & \alpha I_m
\end{pmatrix}.
\]

Следователно

\[
DT(\phi^\ast)-I=
\begin{pmatrix}
\alpha-1 & 0 \\
0 & (\alpha-1) I_m
\end{pmatrix}.
\]

Собствените стойности са равни на $\alpha-1$, поради което REE е E-устойчиво тогава и само тогава, когато

\[
\alpha < 1.
\]

\subsection*{Обобщение}

E-устойчивостта свързва устойчивостта на равновесието под адаптивно обучение с локалните свойства на оператора $T$. Неподвижна точка на $T$ е устойчива, ако малки отклонения на вярванията затихват според асоциираната непрекъсната динамика. В линейния пример това условие се свежда до $\alpha<1$, което съвпада с условието за сходимост на най-малките квадрати.
